
\subsubsection{Research Questions}

The evaluation is structured around four cascading questions:

\begin{enumerate}
\item Does coupling exist? (h-e1) Establishing whether co-occurrence patterns exceed chance levels. Without detectable coupling, subsequent analyses are moot.
\end{enumerate}

\begin{enumerate}
\item Is coupling difficulty-independent? (h-m1) Distinguishing genuine shared vulnerabilities from spurious artifacts where hard instances fail on all dimensions simultaneously.
\end{enumerate}

\begin{enumerate}
\item How broad is coupling? (h-m2) Testing whether coupling spans many dimension pairs (broad architectural bottlenecks) or few pairs (targeted vulnerability clusters).
\end{enumerate}

\begin{enumerate}
\item Are coupling profiles model-specific? (h-c1) Testing whether different models exhibit distinct coupling patterns reflecting different training approaches or architectures.
\end{enumerate}

\subsubsection{Data}

Synthetic coupling data was generated to emulate production benchmark structure. Three model variants (GPT-4, Claude-3-Sonnet, Llama-3-70B) were simulated with 500 instances per model and 100 instances per dimension (truthfulness, robustness, fairness, safety, privacy). Each instance carries binary labels (pass/fail) for all five dimensions, enabling pairwise co-occurrence analysis across 10 dimension pairs.

Coupling patterns were initialized with target correlation values: phi 0.35-0.40 for truthfulness-robustness and fairness-safety pairs, phi 0.15-0.25 for remaining pairs. Model-specific variation was introduced via 7\% noise perturbations. Difficulty scores were simulated as Normal(0.5, 0.15) with independence constraint |correlation(difficulty, dimension)| < 0.2.

This synthetic data validates measurement methodology capabilities but does not characterize real-world coupling patterns. MultiTrust and TrustLLM benchmarks remain access-gated; real LLM API evaluations were not conducted in this proof-of-concept phase. Coupling magnitudes shown demonstrate detectability if coupling exists, not that these specific values exist in production models.

\subsubsection{Experimental Design}

\textbf{h-e1: Coupling Existence Test}

Method: Phi coefficient analysis on $2\times 2$ contingency tables (pass/fail on dimension A vs dimension B).

Threshold: phi $\geq$ 0.3 (medium effect size) and p < 0.01 (statistical significance).

Gate criterion: $\geq 1$ model exhibits $\geq 1$ significant dimension pair.

\textbf{h-m1: Difficulty-Independence Validation}

Method: Dual validation---(1) partial correlation controlling for instance difficulty via partial phi coefficient, (2) quartile stratification computing phi within difficulty quartiles (Q0=easy, Q3=hard).

Threshold: partial phi $\geq$ 0.25 for $\geq 2$ dimension pairs and coupling observable in $\geq 3$ of 4 difficulty quartiles.

Gate criterion: MUST\_WORK. Difficulty-independence distinguishes shared vulnerability mechanisms from measurement artifacts.

\textbf{h-m2: Coupling Breadth Analysis}

Method: Count significant dimension pairs per model with Bonferroni correction (alpha = 0.01/30 $\approx$ 0.00033).

Threshold: $\geq 3$ significant pairs (phi $\geq$ 0.3, p\_adj < 0.01) for $\geq 2$ of 3 models.

Gate criterion: SHOULD\_WORK. This tests sparse vs broad coupling hypothesis.

\textbf{h-c1: Model-Specific Fingerprints}

Method: Mantel test comparing coupling matrices across model pairs via permutation test (10,000 iterations).

Threshold: Mantel r < 0.7 (low similarity) and p < 0.0167 (Bonferroni-corrected for 3 model pairs) for $\geq 1$ model pair.

Gate criterion: SHOULD\_WORK.

\subsubsection{Metrics}

\begin{itemize}
\item \textbf{Phi coefficient:} Effect size for $2\times 2$ contingency tables, computed as sqrt(chi-square / n). Ranges [0,1] with thresholds 0.1=small, 0.3=medium, 0.5=large.
\item \textbf{Chi-square p-value:} Statistical significance test for independence hypothesis. Alpha = 0.01.
\item \textbf{Partial phi:} Phi coefficient after removing linear relationship with difficulty score, implemented via pingouin.partial\_corr.
\item \textbf{Mantel r:} Pearson correlation between vectorized upper-triangle entries of two coupling matrices, with permutation-based p-value.
\end{itemize}

\subsubsection{Statistical Considerations}

Multiple testing correction: h-m2 applies Bonferroni correction (alpha\_adj = 0.01/30). h-e1 and h-m1 test prespecified pairs (truthfulness-robustness, fairness-safety) without correction, justified by prior work documenting these couplings.

Power analysis: Mantel test requires n $\geq$ 500 instances per dimension for detecting r = 0.5 at 80\% power \citep{Legendre2010}. Sample size (n = 100 per dimension) is underpowered for h-c1.

Difficulty independence assumption: Synthetic difficulty scores were constrained to |correlation(difficulty, dimension)| < 0.2, validating independence assumption for partial correlation.
